\documentclass[11pt]{amsart}

\newcommand{\AUTHORNAME}{Yitzchak Shmalo}
\newcommand{\AFFILIATION}{Einstein Institute of Mathematics, The Hebrew University of Jerusalem, Jerusalem, Israel}
\newcommand{\AUTHOREMAIL}{yitzchak.shmalo@gmail.com}
\newcommand{\FUNDING}{The research presented in this paper was supported by the European Research Council (ERC) under the European Union's Horizon Europe research and innovation programme (grant agreement No.~101041711), by the Simons Foundation as part of the Collaboration on the Mathematical and Scientific Foundations of Deep Learning, by Heights Labs, by Convex Nexus Capital, by the Israel Science Foundation (grant number 2258/19), and by the Israel Science Foundation (ISF Grant 4101/25).}
\newcommand{\ACKNOWLEDGMENTS}{The final manuscript was audited by Claude Opus~5.5 (Anthropic), which checked the proofs and the cited results against the sources. }

\usepackage[utf8]{inputenc}
\usepackage[T1]{fontenc}
\usepackage{lmodern}
\usepackage{amsmath,amssymb,amsthm}
\usepackage[margin=1.15in]{geometry}
\usepackage[colorlinks=true,linkcolor=blue,citecolor=blue,urlcolor=blue]{hyperref}
\usepackage[capitalise,nameinlink]{cleveref}

\newtheorem{theorem}{Theorem}[section]
\newtheorem{proposition}[theorem]{Proposition}
\newtheorem{lemma}[theorem]{Lemma}
\newtheorem{corollary}[theorem]{Corollary}
\theoremstyle{remark}
\newtheorem{remark}[theorem]{Remark}

\crefname{theorem}{Theorem}{Theorems}
\crefname{proposition}{Proposition}{Propositions}
\crefname{lemma}{Lemma}{Lemmas}
\crefname{corollary}{Corollary}{Corollaries}
\crefname{remark}{Remark}{Remarks}
\crefname{section}{Section}{Sections}

\DeclareMathOperator{\br}{br}
\DeclareMathOperator{\dist}{dist}

\begin{document}

\title[Branching number of geodesic spanning trees]{Every geodesic spanning tree of a Cayley graph has branching number equal to the growth rate}

\author{\AUTHORNAME}
\address{\AFFILIATION}
\email{\AUTHOREMAIL}

\subjclass[2020]{Primary 05C05, 60K35; Secondary 20F65, 05C25}
\keywords{Branching number, growth rate, geodesic spanning tree, Cayley graph, quasi-transitive graph, percolation on trees}

\begin{abstract}
A spanning tree $T$ of a graph $G$, rooted at a vertex $o$, is geodesic if the distance from $o$ to every vertex is the same in $T$ as in $G$. Lyons and Peres asked whether every geodesic spanning tree of the Cayley graph of a finitely generated group of growth rate $b$ has branching number $b$; they proved this for the lexicographically minimal tree. We show that the answer is yes, and more generally that every geodesic spanning tree of an infinite, connected, locally finite, quasi-transitive graph has branching number equal to the growth rate of the graph. Consequently the critical probability of Bernoulli percolation on any such tree is $1/b$.
\end{abstract}

\maketitle

\section{Introduction}

Let $G$ be an infinite, connected, locally finite graph and $o$ a vertex of $G$. Write $|x|=\dist_G(o,x)$, $B_n(x)$ for the ball of radius $n$ about $x$ and $S_n=\{x:|x|=n\}$. A spanning tree $T$ of $G$ rooted at $o$ is a \emph{geodesic subtree} if $\dist_T(o,x)=|x|$ for every vertex $x$. Such trees exist in every connected graph: give each $x\ne o$ a parent among its neighbours in $S_{|x|-1}$ (a breadth-first search tree is an example). Level $n$ of a geodesic spanning tree is exactly the sphere $S_n$.

The branching number of a tree $T$ rooted at $o$ is
\[
\br T=\sup\Bigl\{\lambda\ge1:\ \inf_{\Pi}\sum_{e\in\Pi}\lambda^{-|e|}>0\Bigr\},
\]
where $\Pi$ ranges over the finite sets of edges separating $o$ from infinity and $|e|$ is the distance from $o$ to the endpoint of $e$ farther from $o$ \cite[Section~1.2]{LP}. By a theorem of Lyons \cite{Lyons1990}, $\br T$ governs random walk and percolation on $T$; in particular the critical probability of Bernoulli percolation on $T$ is $1/\br T$.

Lyons and Peres \cite[Question~3.12]{LP} ask the following. \emph{If $T$ is a geodesic spanning tree of the Cayley graph of a finitely generated group of growth rate $b$, is $\br T=b$?} They show that the answer is yes for the lexicographically minimal geodesic tree, which is subperiodic, so that its branching number equals its growth \cite[Section~3.4]{LP}. A general geodesic tree need not be subperiodic: the choice of parent at each vertex is arbitrary, and the subtrees below different vertices of the same sphere can look very different. Liu and Xiang \cite[Section~3.2]{LX} obtain $\br T=\operatorname{gr}G$ in a quasi-transitive graph after constructing a lexicographic geodesic tree. Their Lemma~3.5 is worded more generally, but its proof goes through subperiodicity and does not cover an arbitrary geodesic spanning tree. We give a direct proof for every such tree.

\begin{theorem}\label{thm:main}
Let $G$ be an infinite, connected, locally finite, quasi-transitive graph with growth rate $b$, and let $T$ be a geodesic spanning tree of $G$ rooted at $o$. Then $\br T=b$. If $G$ is vertex-transitive, then $\sum_{e\in\Pi}b^{-|e|}\ge 1$ for every finite edge cutset $\Pi$ of $T$.
\end{theorem}

\begin{corollary}\label[corollary]{cor:cayley}
Every geodesic spanning tree of a Cayley graph of an infinite finitely generated group of growth rate $b$ has branching number $b$, and Bernoulli($p$) percolation on it has critical probability $1/b$.
\end{corollary}

The proof uses no structure of $T$ beyond the geodesic property. Seen from its own root, the subtree below a vertex $v$ embeds into $G$ preserving distances from its root, and transitivity lets us compare every such subtree with the whole graph at once. The comparison is made with a generating function evaluated just above $b$, where it blows up.

Some transitivity is needed, because level sizes alone do not determine the branching number. Hang a finite binary tree of height $n$ from the $n$th vertex of a ray. The levels of the resulting tree grow like $2^{n/2}$, but a single edge of the ray at depth $N$ is a cutset, so the branching number is $1$.

\section{Proof}

Throughout, $G$ is as in \cref{thm:main}. Fix representatives $o_1,\dots,o_k$ of the orbits of the automorphism group of $G$ and set $R=\max_i\dist(o,o_i)$. Since $G$ is locally finite and connected, $|B_n(x)|$ depends only on the orbit of $x$.

\begin{lemma}\label[lemma]{lem:growth}
Let $f(n)=\max_i|B_n(o_i)|$. The limit $b=\lim_n f(n)^{1/n}$ exists, equals $\lim_n|B_n(o)|^{1/n}$, and $|B_n(o)|\ge b^{\,n-R}$ for every $n\ge R$.
\end{lemma}

\begin{proof}
Every point of $B_{m+n}(o_j)$ lies within distance $n$ of a point of $B_m(o_j)$, so $B_{m+n}(o_j)$ is covered by the balls $B_n(x)$, $x\in B_m(o_j)$, each of size at most $f(n)$. Hence $f(m+n)\le f(m)f(n)$, and by Fekete's lemma $b=\lim f(n)^{1/n}=\inf_n f(n)^{1/n}$ exists and $f(n)\ge b^n$ for all $n$. Since $B_{n-R}(o_i)\subseteq B_n(o)\subseteq B_{n+R}(o_i)$ for each $i$, we get $f(n-R)\le|B_n(o)|\le f(n+R)$, which gives both remaining claims.
\end{proof}

For $x\in V(G)$ and $t>b$ put $Z_x(t)=\sum_{y}t^{-\dist(x,y)}$ and $Z(t)=Z_o(t)$. Summing by spheres and then by parts,
\begin{equation}\label{eq:Z}
Z(t)=\sum_{n\ge0}|S_n|\,t^{-n}=\Bigl(1-\frac1t\Bigr)\sum_{n\ge0}|B_n(o)|\,t^{-n},
\end{equation}
which is finite for $t>b$ by \cref{lem:growth}.

\begin{lemma}\label[lemma]{lem:blowup}
If $b>1$, then $Z(t)\to\infty$ as $t\downarrow b$.
\end{lemma}

\begin{proof}
By \eqref{eq:Z} and \cref{lem:growth}, $Z(t)\ge(1-1/t)\,b^{-R}\sum_{n\ge R}(b/t)^n$, and the series diverges as $t\downarrow b$.
\end{proof}

For a vertex $v$ write $D(v)$ for the set consisting of $v$ and its descendants in $T$.

\begin{lemma}\label[lemma]{lem:radial}
If $x\in D(v)$, then $\dist_G(v,x)=|x|-|v|$.
\end{lemma}

\begin{proof}
The path from $o$ to $x$ in $T$ passes through $v$, and $T$ is geodesic, so $|x|=|v|+\dist_T(v,x)\ge|v|+\dist_G(v,x)\ge|x|$.
\end{proof}

\begin{lemma}\label[lemma]{lem:subtree}
For every vertex $v$ and every $t>b$,
\[
\sum_{x\in D(v)}t^{-|x|}\le t^{R}\,t^{-|v|}\,Z(t).
\]
\end{lemma}

\begin{proof}
Choose $i$ and an automorphism $\varphi$ of $G$ with $\varphi(v)=o_i$. By \cref{lem:radial}, $|x|=|v|+\dist(o_i,\varphi(x))$ for $x\in D(v)$, and $\varphi$ is injective, so
\[
\sum_{x\in D(v)}t^{-|x|}=t^{-|v|}\sum_{x\in D(v)}t^{-\dist(o_i,\varphi(x))}\le t^{-|v|}Z_{o_i}(t).
\]
Finally $\dist(o,y)\le R+\dist(o_i,y)$ for every $y$, so $Z_{o_i}(t)\le t^{R}Z(t)$, since $t>b\ge1$.
\end{proof}

\begin{proof}[Proof of \cref{thm:main}]
\emph{Upper bound.} The edges at level $n$ of $T$ form a cutset of size $|S_n|\le|B_n(o)|$. For $\lambda>b$, \cref{lem:growth} gives $\sum_{|e|=n}\lambda^{-n}\le|B_n(o)|\lambda^{-n}\to0$, so $\br T\le b$.

\emph{Lower bound.} Every finite edge cutset has at least one edge, so $\br T\ge1$ always, and we may assume $b>1$. Let $\Pi$ be a finite edge cutset separating $o$ from infinity, let $P$ be the set of far endpoints of the edges of $\Pi$, and let $K$ be the component of $o$ in $T$ with the edges of $\Pi$ removed. If $K$ were infinite, K\H{o}nig's lemma would give an infinite ray from $o$ inside $K$, which avoids $\Pi$; so $K$ is finite. If $x\notin K$, the path from $o$ to $x$ in $T$ leaves $K$ through an edge of $\Pi$, whose far endpoint $v\in P$ is an ancestor of $x$; thus $x\in D(v)$. Using \cref{lem:subtree}, for $t>b$,
\[
Z(t)=\sum_{x}t^{-|x|}\le\sum_{x\in K}t^{-|x|}+\sum_{v\in P}\ \sum_{x\in D(v)}t^{-|x|}\le\sum_{x\in K}t^{-|x|}+t^{R}Z(t)\sum_{e\in\Pi}t^{-|e|}.
\]
Divide by $Z(t)$ and let $t\downarrow b$. The first term on the right is at most $|K|$, so by \cref{lem:blowup} it vanishes after division, and we obtain
\[
\sum_{e\in\Pi}b^{-|e|}\ge b^{-R}.
\]
The bound is uniform in $\Pi$, so $\br T\ge b$. In the vertex-transitive case take $o_1=o$; then $R=0$ and the bound reads $\sum_{e\in\Pi}b^{-|e|}\ge1$.
\end{proof}

\begin{proof}[Proof of \cref{cor:cayley}]
The question concerns infinite groups. A Cayley graph of an infinite finitely generated group with respect to a finite generating set is infinite, connected, locally finite and vertex-transitive, so \cref{thm:main} applies. The statement on percolation is Lyons's theorem $p_c(T)=1/\br T$ \cite{Lyons1990}.
\end{proof}

\begin{remark}
Dead ends of $G$, such as those of the lamplighter group \cite[Section~3.4]{LP}, cause no difficulty: a subtree $D(v)$ may be finite, and \cref{lem:subtree} only bounds it from above. The proof never asks an individual subtree to be large; the lower bound comes from the whole graph through $Z(t)$, and for a fixed cutset only the finitely many vertices of $K$ are lost.
\end{remark}

\section*{Acknowledgements and funding}
\ACKNOWLEDGMENTS\FUNDING

\section*{Declaration of competing interest}
The author holds equity in Heights Labs and in Convex Nexus Capital and has a management role at Convex Nexus Capital. These entities are included in the funding declaration above.

\section*{Declaration of generative AI and AI-assisted technologies}
The paper comes out of Math Brain, an ongoing project on the use of AI systems for research mathematics. In October 2026 an OpenAI GPT model, running as a research task of that project, proposed the proof of \cref{thm:main} for Cayley graphs. Claude Opus~5.5 (Anthropic), working in Claude Code, audited the proof against the book, extended it to quasi-transitive graphs and wrote the paper, and a second, independent model audit checked it. These checks are not peer review. The author directed the work and is responsible for the content of the paper and for decisions concerning its dissemination.

\bibliographystyle{amsplain}
\bibliography{references}

\end{document}
